% before: lit-open-inputs
Most results in this chapter are proved outright. A few concern the local
Langlands correspondence and Galois representations attached to automorphic
representations, and those are proved from the two theorems of
Chapter~\ref{chapter-galois-representations}. The remark says which results
depend on them and why the dependence is harmless.

% before: lit-bundle-unitary
The definite side of the argument works on unitary groups over a CM field and
passes between automorphic representations, their local components and
Galois representations. It draws on representation theory of real and
$p$-adic groups, Eisenstein series, base change and class field theory. The
theorem lists the results it uses with their published sources.

% before: lit-bundle-theta
The theta lifts of Chapter~\ref{chapter-unitary-theta} need a second group of
inputs: models of the Weil representation, local properties of discrete
series and Hermitian forms, and analytic facts about cusp forms and Hecke
$L$-functions. Several entries repeat entries of the previous list, since the
two lists feed different parts of the proof.

% before: lit-casselman-wallach
Automorphic forms on unitary groups are compared through their smooth vectors
at the archimedean places, while the algebraic arguments see only
$K$-finite vectors. The globalization theorem lets a map defined on
$K$-finite vectors extend, continuously and uniquely, to all smooth vectors.

% before: lit-howe-duality-real
At the archimedean places the theta correspondence is studied on the Fock
space. The lifts used later start from the trivial representation of the
smaller group, so only this case of Howe's theorem is needed: the big theta
lift of the trivial representation has finite length and at most one
irreducible quotient.

% before: lit-green-cuspidal
The lower bound for the slope of the glued lattice uses cuspidal characters
of general linear groups over finite fields. Green's theorem
parametrises these characters by regular characters of the multiplicative
group of an extension field and gives their values on elliptic regular
elements.

% before: lit-labelling
To compare cuspidal characters for different degrees, the argument names them
by irreducible polynomials. This needs a consistent choice of roots of unity
and of generators of the multiplicative groups of the finite fields.

% before: lit-green-labelled
With a labelling fixed, the characters of Green's theorem can be indexed by
polynomials. The covolume bound and the assembly of the proof use the
statement in this form.

% before: lit-weil-datum
The global theta correspondence is realised on functions on an adelic
symplectic space that are quasi-periodic under its rational points. Evaluation
at the origin then gives the theta kernel, and the Siegel-Weil formula
compares its integral with an Eisenstein series.

% before: lit-siegel-eisenstein
The Eisenstein series in the Siegel-Weil formula and in the doubling method
live on the unitary group of a doubled space. They are built from sections of
a degenerate principal series whose restriction to a fixed maximal compact
subgroup does not depend on $s$, so that only the character on the Siegel
parabolic varies.

% before: lit-weil-siegel
The near-orthogonality of the columns of the lattice of lifted forms rests on
integrals of theta kernels. Weil's theorem says when these integrals converge,
and in the range of absolute convergence of the Eisenstein series identifies
them with an Eisenstein series.

% before: lit-siegel-weil
Some of the theta integrals used lie outside the range in which the Eisenstein
series converges, though the theta integral itself still converges. In that
range the identity continues to hold for Fock vectors, with the Eisenstein
series evaluated by analytic continuation.

% before: lit-doubling
The Rallis inner product formula computes norms of theta lifts through the
doubling method. It needs the continuation of Siegel Eisenstein series and
the unfolding of the doubling integral against two cusp forms.

% before: lit-vacuous-slot
One entry of the list of Theorem~\ref{lit-bundle-unitary} has no
mathematical content. Since its name refers to a real theorem, the remark says
what it contains and what follows for the doubling method.

% before: lit-jacquet-shalika
Automorphic representations of $GL_n$ over CM fields are identified by their
Satake parameters at almost all places. The Jacquet-Shalika bound on those
parameters keeps the local Euler factors under control.

% before: lit-clozel-cohomological
The Galois representations of Chapter~\ref{chapter-galois-representations}
are attached to cohomological cuspidal representations. Clozel's lemma shows
that regular algebraic cuspidal representations are of this kind.

% before: lit-bhatt-scholze
The Galois-theoretic arguments work with lattices in crystalline
representations of local Galois groups. Bhatt and Scholze describe them through
prismatic $F$-crystals, and the full faithfulness of their realization
transfers maps between lattices to maps between crystals.

% before: lit-hodge-tate-sequence
The arithmetic volume estimate of Chapter~\ref{chapter-covolumes} works with
Hodge bundles of abelian schemes and the Tate modules of their $p$-divisible
groups. The duality between the
Hodge-Tate maps of a $p$-divisible group and of its dual is the part of
Fargues's theorem that this comparison needs.

% before: lit-linnik
Some local choices on the theta side need a prime in a prescribed residue
class. Linnik's theorem gives one whose size is bounded by a power of the
modulus.

% before: lit-chebotarev-degree-one
Galois-theoretic arguments on the definite side need places with a prescribed
Frobenius that avoid a finite set of bad places. Places of degree one suffice,
and the qualitative Chebotarev theorem provides them.

% before: lit-snc-sequence
Principalization produces a sequence of blow-ups whose centres meet the
existing boundary transversally. The definition records the conditions imposed
at each step.

% before: lit-principalization
The arithmetic volume estimate of Chapter~\ref{chapter-covolumes} needs to
make ideal sheaves on smooth varieties invertible by blow-ups in smooth
centres. The boundary already present must stay in normal crossings with each
centre.

% before: lit-principalization-projective
For the arithmetic Hilbert-Samuel formula the ideal lives on a smooth
projective variety, and the blown-up variety must again be smooth and
projective of the same dimension. The corollary specializes the theorem to
this case.

% before: lit-resolution
The volume estimate also needs a smooth model of a singular variety with a
normal crossings boundary over its singular locus. This is the form of
resolution of singularities it uses.

% before: lit-arithmetic-variety
Arakelov theory measures sections of line bundles on models over rings of
integers by their norms at the complex places. The definition sets up the
objects of the next two statements.

% before: lit-arithmetic-hilbert-samuel
Covolumes of lattices of sections are counted through small sections. For a
nef Hermitian line bundle, the arithmetic Hilbert-Samuel formula gives the
growth of their number in terms of the arithmetic self-intersection.

% before: lit-multiprojective
The covolume bound of Chapter~\ref{chapter-covolumes} counts small sections on
subvarieties of products of projective spaces over $\mathbf Z$, given by
multihomogeneous prime ideals. Sections are written as quotients of
multihomogeneous polynomials, and arithmetic intersection numbers enter only
through the inductive formula along divisors of sections, which is part of
the definition.

% before: lit-hilbert-samuel-upper
The covolume bound needs only an upper bound for the number of small
sections, together with nonnegativity of self-intersections of nef bundles.
Two forms of the upper bound are used: one for bundles generated by finitely
many small sections, and one for bundles whose monomial sections are small,
with sections counted on the normalization.

% before: lit-zarhin
The estimates on abelian schemes are easier for principally polarized ones.
Zarhin's trick reduces a polarized abelian scheme to a direct factor of a
principally polarized one.
